\section{Future Work}
\label{sec:future}

Future work is tiered by operational maturity, following the
companion-paper convention \cite{blackrim_moa_paper}: Tier~1
items are committed deliverables for the next deployment cycle;
Tier~2 items are refinements that depend on accumulated telemetry;
Tier~3 items are exploratory and may not ship.

\subsection{Tier 1 --- next deployment cycle}
\label{sec:future:tier1}

\begin{enumerate}[leftmargin=*]
  \item \textbf{Pattern-DSL implementation
        (\cref{sec:system:patterndsl}).} Two backends: optional
        shell-out to the \texttt{ast-grep} binary when installed;
        in-tree Go fallback for the subset of the pattern grammar
        the registered intents exercise. Authoring-velocity target:
        $\leq 50$~LoC of YAML per new tier-0 intent, down from
        $\sim 200$~LoC of Go today. Gates OQ-AST-6.
  \item \textbf{Outline telemetry fill-in.} Run the discipline
        subsystem in warn-with-auto-prepend mode for 14~days across
        the Blackrim production agent fleet; accumulate
        \texttt{outline-events.jsonl} to the volume required for
        OQ-AST-3 (adoption trajectory) and OQ-AST-5 (end-to-end
        token reduction on paired traces).
  \item \textbf{Polyglot bench corpus expansion.} Extend
        \texttt{cmd/gt/testdata/bench/} from the current
        four-file fixture to a $\geq 50$-file stratified corpus
        per language. Gates OQ-AST-1 (token-budget conformance).
\end{enumerate}

\subsection{Tier 2 --- post-deployment refinement}
\label{sec:future:tier2}

\begin{enumerate}[leftmargin=*]
  \item \textbf{Grammar-version abstraction layer.} A
        language-neutral category table that maps grammar
        node-type names to canonical categories
        (\texttt{function-def}, \texttt{class-def},
        \texttt{exception-handler}, etc.). Reduces the
        per-language-grammar maintenance surface to a per-language
        category map plus a single shared significance rule over
        categories.
        Addresses \cref{sec:discussion:treesitter}.
  \item \textbf{Outline caching.} Currently every \texttt{Read}
        on a $\loc \geq \threshold$ file re-parses the file. If
        OQ-AST-2 latency turns out to be a user-perceived
        bottleneck, a content-addressed outline cache keyed on
        $(\text{file mtime}, \text{file size}, \text{language})$
        can be added without changing the emitter contract.
  \item \textbf{Cross-file outline composition.} The
        \texttt{codeindex\_explore} surface already returns
        \texttt{source\_sections} and \texttt{additional\_files};
        a \emph{cross-file outline} that summarises the
        relationship between a set of files (call graph, type
        dependency, package structure) is the natural extension.
        Differs from per-file outline in being a graph view, not a
        tree view.
  \item \textbf{Semantic-fidelity gate.} The current
        \cref{alg:gate} runs native syntax check and (where
        available) LSP-diagnose. A semantic-fidelity gate would
        additionally run a small test subset (the package's unit
        tests, with a wall-clock budget) on the candidate diff
        before accepting. Costs latency, buys depth.
\end{enumerate}

\subsection{Tier 3 --- parking lot}
\label{sec:future:tier3}

\begin{enumerate}[leftmargin=*]
  \item \textbf{Outline-as-context for non-code artefacts.}
        Markdown, LaTeX, and structured-config files
        (Kubernetes manifests, Terraform modules) have AST-like
        structure that the same lens could expose. The
        cost-benefit is materially different (Markdown is already
        outline-shaped) so the question is whether the unified
        emitter or separate emitters per artefact class.
  \item \textbf{Outline diffing.} An outline-diff between
        two revisions of a file shows what \emph{structurally}
        changed (declarations added / removed, public API
        delta) independent of what cosmetically changed. Useful
        for code review and for change-classification (is this PR
        a refactor or a behaviour change?). The plan-token
        mechanism in \cref{sec:algorithm:planexecute} is a
        first step.
  \item \textbf{Adversarial robustness of sanitisation.}
        Empirical red-teaming of the sanitisation pass: can an
        adversary craft package-doc content that survives the
        sanitiser but successfully prompt-injects an outline
        consumer? If so, what is the minimal escape pattern set
        that closes those holes?
  \item \textbf{Federated outline registry.} For multi-repo
        agent workflows: a registry of outlines computed at
        clone-time, served from a CDN. Trades freshness for
        latency. Differs from the codeindex symbol-graph in
        emitting outlines (~$300$~tokens each) rather than
        symbols.
\end{enumerate}

\subsection{Pre-registered open questions}

The full list of pre-registered open questions
\texttt{OQ-AST-1..OQ-AST-7} is in \cref{app:open-questions}. They
are the empirical questions whose answers gate promotion or
demotion of the various subsystems. Five of them
(\texttt{OQ-AST-1..OQ-AST-5}) are the pending claims of
\cref{sec:eval}; the remaining two are forward-looking
(\texttt{OQ-AST-6}: pattern-DSL fidelity vs. native parser walking;
\texttt{OQ-AST-7}: cross-file outline composition utility on
multi-file refactors).
